\documentclass[12pt]{article}
\usepackage[utf8]{inputenc}
\usepackage[T1]{fontenc}
\usepackage{hyperref}
\title{Punan Aput Language Resources for K-12 Curriculum Development}
\author{Compiled by Claude Sonnet 5.5 (claude-sonnet-5-5) for Andr\'as Kornai}
\date{}
\begin{document}
\maketitle

\section{Introduction}
Punan Aput (ISO 639-3 \texttt{pud}, Glottocode puna1267) is an Austronesian language of Kalimantan, Indonesia. Sources disagree on its location (East, West or North Kalimantan; one data sheet places it in Malinau regency, south of the Kayan river) and on speaker numbers (about 370 in a 1981 count cited by Wikipedia; about 700 in Joshua Project, citing Ethnologue 2016). Very little material exists. In searches I found no primer, textbook, dictionary, grammar, orthography or scripture for this language, and no school programme. This report lists only general reference pages.

\section{Primary Educational Resources}
None found. No primers, textbooks or school materials in Punan Aput turned up in searches.

\section{Linguistic and Reference Materials}
\begin{itemize}
\item Aput language, Wikipedia (English): short overview of classification, location and speaker estimates. \url{https://en.wikipedia.org/wiki/Aput_language}
\item Punan Aput, Kalimantan: language data sheet of the numerals project (MPI-EVA), a numeral-system data page. \url{https://lingweb.eva.mpg.de/channumerals/Punan-Aput.htm}
\item Soriente and Inagaki, 2012 overview slides on Kalimantan languages and documentation; notes that the classification of the Punan languages is unresolved and that documentation has been limited. \url{https://lingdy.aa-ken.jp/wp-content/uploads/2012/01/120217_soriente_inagaki_h2.pdf}
\end{itemize}

\section{Religious and Community Resources}
\begin{itemize}
\item Aput Punan in Indonesia, Joshua Project: people-group profile (population, language estimates). \url{https://joshuaproject.net/people_groups/14475}
\end{itemize}
No community organisation or scripture translation was found.

\section{Digital and Audio Resources}
None found.

\section{Curriculum Development Context}
The language is reported as endangered, with few speakers and almost no published documentation. Any curriculum would first require community-led documentation (word lists, texts, recordings) and an agreed orthography, working with Aput Punan speakers. Indonesian would remain the language of schooling in the meantime. Verify all details with the sources above, which are general web pages and are not themselves teaching materials.

\end{document}